% Generated draft: requires mathematical and source review.
\documentclass[11pt]{article}
\usepackage[margin=25mm]{geometry}
\usepackage{fontspec}
\setmainfont{Latin Modern Roman}
\usepackage{amsmath,amssymb,mathtools,unicode-math}
\setmathfont{Latin Modern Math}
\usepackage{enumitem,needspace,xurl,hyperref}
\setlength{\parindent}{0pt}
\setlength{\parskip}{5pt}
\setlength{\emergencystretch}{3em}
\newcommand{\N}{\mathbb{N}}
\newcommand{\Z}{\mathbb{Z}}
\newcommand{\Q}{\mathbb{Q}}
\newcommand{\R}{\mathbb{R}}
\newcommand{\C}{\mathbb{C}}
\newcommand{\F}{\mathbb{F}}
\newcommand{\A}{\mathbb{A}}
\newcommand{\PP}{\mathbb{P}}
\newcommand{\E}{\mathbb{E}}
\newcommand{\eps}{\varepsilon}
\newcommand{\dd}{\,\mathrm{d}}
\begin{document}
\section*{UnsolvedMath ID: 30001579}
\textbf{Source problem code:} catalogue entry
\subsection{Definitions and mathematical statement}
Let \(\Gamma\) be a finitely generated residually finite group. Here residually finite means that for every \(g\in\Gamma\setminus\{1\}\), there is a homomorphism from \(\Gamma\) to a finite group whose value at \(g\) is not the identity. For a finitely generated group \(H\), let \(d(H)\) denote the minimum cardinality of a generating set of \(H\).

Interpret a normal chain with trivial intersection in the standard sense as a sequence
\[
\Gamma\geq\Gamma_1\geq\Gamma_2\geq\cdots,
\qquad \Gamma_n\trianglelefteq\Gamma,
\qquad [\Gamma:\Gamma_n]<\infty,
\qquad \bigcap_{n\geq1}\Gamma_n=\{1\}.
\]
The finite-index requirement is implicit in the supplied rank-gradient formulation; it is made explicit here. No assumption of strict inclusions, torsion-freeness, or infinitude of \(\Gamma\) is imposed. Define
\[
\operatorname{RG}(\Gamma,(\Gamma_n))
=\lim_{n\to\infty}\frac{d(\Gamma_n)-1}{[\Gamma:\Gamma_n]}.
\]
For such chains this limit exists by the Schreier index inequality.

\textbf{Conjecture (chain independence of rank gradient).} For every finitely generated residually finite group \(\Gamma\) and every two descending chains \((\Gamma_n)_{n\geq1}\) and \((\Lambda_n)_{n\geq1}\) of finite-index subgroups normal in \(\Gamma\), each with trivial intersection, one has
\[
\operatorname{RG}(\Gamma,(\Gamma_n))
=\operatorname{RG}(\Gamma,(\Lambda_n)).
\]
This is the conjecture recorded by the supplied source, not a claimed theorem. The source's label that it is open is supplied dated curation, not an independently verified status assessment.
\subsection{Short English statement}
For a finitely generated residually finite group, consider descending chains of finite-index normal subgroups whose intersection is the identity. The conjecture asks whether the limiting number of generators minus one, divided by subgroup index, has the same value for every such chain.
\subsection{Sources}
\begin{itemize}
\item Local UnsolvedMath registry, \texttt{data/unsolvedmath/problems.json}, permanent ID \(30001579\), title ``Chain Independence of Rank Gradient,'' exact source code \texttt{catalogue entry}, in the supplied collection ``Oberwolfach Reports Open Problems.'' No individual proposer is identified in the supplied record; its proposed year is 2010.
\item Supplied report citation: \emph{Actions and Invariants of Residually Finite Groups: Asymptotic Methods} (2011), Oberwolfach Reports, pp. 2335--2391, DOI 10.4171/owr/2010/40. \url{https://doi.org/10.4171/owr/2010/40}. The supplied background describes these reports as documenting mathematical workshops and their open-problem discussions. The supplied literature-source label is ``original OWR report.''
\item Supplied literature reference: ``Abert--Nikolov, rank gradient and cost.'' \url{https://ems.press/journals/jems/articles/4244}.
\item Literature checked at: \texttt{2026-08-21}. This is supplied dated curation, not independently verified. Its assessment labels the problem open and states that chain independence remains unresolved in general, while results relating rank gradient to cost for Farber chains do not establish independence for all normal residual chains. The supplied background explicitly requests independent verification of time-sensitive status and cited claims.
\item The registry labels the statement \texttt{corrected\_verified} and reports that its mathematical content was checked against the official Oberwolfach report, with notation, context, and question form normalized without changing the source-level claim. These are supplied verification claims; no source visit or independent verification is asserted here.
\end{itemize}
\end{document}
